% =============================================================================
% DISSERTAÇÃO/TESE MODELO — ABNT (abnTeX2)
% =============================================================================
% Classe: abntex2.cls v1.9.7
% Normas: NBR 14724 (estrutura), NBR 10520:2023 (citações),
%         NBR 6023:2018 (referências), NBR 6024 (numeração),
%         NBR 6027 (sumário), NBR 6028:2021 (resumo)
% =============================================================================
\documentclass[
  12pt,
  openright,
  twoside,
  a4paper,
  brazilian
]{abntex2}

% --- Pacotes essenciais ---
\usepackage[utf8]{inputenc}
\usepackage[T1]{fontenc}
\usepackage{amsmath,amssymb}
\usepackage{graphicx}
\usepackage[brazilian,portuges]{babel}
\usepackage{hyperref}
\usepackage{booktabs}
\usepackage{float}
\usepackage[alf]{abntex2cite}
\usepackage[algoruled,portuguese]{algorithm2e}
\usepackage{multirow}
\usepackage{tabularx}
\usepackage{microtype}
\usepackage{lmodern}
\usepackage{xcolor}

% --- Configuração das margens (NBR 14724) ---
% 3cm superior/esquerda, 2cm inferior/direita
\setlrmarginsandblock{3cm}{2cm}{*}
\setulmarginsandblock{3cm}{2cm}{*}
\checkoddpage
\checkopenright
\checkandfixthelayout

% --- Metadados do PDF ---
\hypersetup{
  pdfauthor={Nome do Autor},
  pdftitle={Título da Dissertação},
  pdfsubject={Área de Concentração},
  colorlinks=true,
  linkcolor=black,
  citecolor=black,
  urlcolor=black
}

% --- Dados institucionais ---
\instituicao{Universidade Federal do Ceará (UFC)}

\begin{document}

% ============================================================================
% ELEMENTOS PRÉ-TEXTUAIS
% ============================================================================

% --- Capa ---
\titulo{Título da Dissertação: Subtítulo (se houver)}
\autor{Nome Completo do Autor}
\local{Fortaleza, CE}
\data{2026}
\orientador{Prof. Dr. Nome do Orientador}
\coorientador{Prof. Dr. Nome do Coorientador (opcional)}
\preambulo{Dissertação apresentada ao Programa de
Pós-Graduação em Tecnologia Educacional da Universidade Federal do Ceará (UFC),
como requisito parcial para a obtenção do título de Mestre em Tecnologia
Educacional.}

\imprimircapa

% --- Folha de rosto ---
\imprimirfolhaderosto

% --- Ficha catalográfica ---
% Deve ser gerada pela biblioteca da instituição e inserida como PDF
% \include{ficha_catalografica}

% --- Folha de aprovação ---
% Deve ser inserida após a defesa. Modelo:
% \begin{folhadeaprovacao}
%   \begin{center}
%     {\textbf{\imprimirautor}}\par
%     \vspace*{\fill}
%     \begin{center}
%       \textbf{\imprimirtitulo}
%     \end{center}
%     \vspace*{\fill}
%     \hspace{.45\textwidth}
%     \begin{minipage}{.5\textwidth}
%       \imprimirpreambulo
%     \end{minipage}%
%     \vspace*{\fill}
%   \end{center}
%   Trabalho aprovado. \imprimirlocal, data da defesa:
%   \assinatura{\textbf{\imprimirorientador} \\ Orientador}
%   \assinatura{\textbf{Convidado 1} \\ Convidado}
%   \assinatura{\textbf{Convidado 2} \\ Convidado}
%   \begin{center}
%     \vfill
%     {\imprimirlocal}\par
%     \imprimirdata
%   \end{center}
% \end{folhadeaprovacao}

% --- Dedicatória (opcional) ---
\begin{dedicatoria}
  \vspace*{\fill}
  \centering
  \noindent
  \textit{Dedico este trabalho a...}
  \vspace*{\fill}
\end{dedicatoria}

% --- Agradecimentos (opcional) ---
\begin{agradecimentos}
Agradeço ao meu orientador, Prof. Dr. Nome, pela orientação e apoio durante
toda a pesquisa.

Aos professores do Programa de Pós-Graduação em Tecnologia Educacional da UFC,
pelas contribuições teóricas e metodológicas.

À [Agência de Fomento], pelo apoio financeiro (Processo nº XXXX).

Aos colegas de laboratório e grupo de pesquisa, pelo diálogo e colaboração.

À minha família, pelo suporte incondicional.
\end{agradecimentos}

% --- Epígrafe (opcional) ---
\begin{epigrafe}
  \vspace*{\fill}
  \begin{flushright}
    \textit{``A citação que inspira o trabalho.''} \\
    (Autor, Ano)
  \end{flushright}
  \vspace*{\fill}
\end{epigrafe}

% --- Resumo (português) ---
\setlength{\absparsep}{18pt}
\begin{resumo}
O resumo deve apresentar objetivo, método, resultados e conclusões da
pesquisa, em parágrafo único de 150 a 500 palavras (dissertação/tese).
Palavras-chave representativas do conteúdo, separadas por ponto e finalizadas
por ponto.
\vspace{\onelineskip}

\noindent
\textbf{Palavras-chave}: Palavra-chave 1. Palavra-chave 2. Palavra-chave 3.
Palavra-chave 4.
\end{resumo}

% --- Abstract (inglês) ---
\begin{resumo}[Abstract]
  \begin{otherlanguage*}{english}
  The abstract should present objective, method, results and conclusions of
  the research in a single paragraph of 150 to 500 words.
  \vspace{\onelineskip}

  \noindent
  \textbf{Keywords}: Keyword 1. Keyword 2. Keyword 3. Keyword 4.
  \end{otherlanguage*}
\end{resumo}

% --- Resumo em outras línguas (opcional) ---
% \begin{resumo}[Résumé]
%   \begin{otherlanguage*}{french}
%   ... texto em francês ...
%   \end{otherlanguage*}
% \end{resumo}

% --- Lista de ilustrações ---
\pdfbookmark[0]{\listfigurename}{lof}
\listoffigures*
\cleardoublepage

% --- Lista de tabelas ---
\pdfbookmark[0]{\listtablename}{lot}
\listoftables*
\cleardoublepage

% --- Lista de abreviaturas e siglas ---
\begin{siglas}
  \item[ABNT] Associação Brasileira de Normas Técnicas
  \item[CAPES] Coordenação de Aperfeiçoamento de Pessoal de Nível Superior
  \item[CNPq] Conselho Nacional de Desenvolvimento Científico e Tecnológico
  \item[UFC] Universidade Federal do Ceará
  \item[TIC] Tecnologias da Informação e Comunicação
\end{siglas}

% --- Lista de símbolos ---
\begin{simbolos}
  \item[$n$] Número de elementos da amostra
  \item[$p$] Nível de significância (p-valor)
  \item[$\sigma$] Desvio-padrão
  \item[$\bar{x}$] Média aritmética
\end{simbolos}

% --- Sumário ---
\pdfbookmark[0]{\contentsname}{toc}
\tableofcontents*
\cleardoublepage

% ============================================================================
% ELEMENTOS TEXTUAIS
% ============================================================================
\textual

% --- 1. INTRODUÇÃO ---
\chapter[Introdução]{Introdução}

A introdução deve apresentar:
\begin{itemize}
  \item Contexto e delimitação do tema
  \item Problema de pesquisa (em forma de pergunta)
  \item Hipóteses ou questões norteadoras
  \item Objetivos (geral e específicos)
  \item Justificativa (relevância teórica, prática e social)
  \item Estrutura do trabalho
\end{itemize}

\section{Problema de Pesquisa}

Formulação clara e delimitada do problema que se pretende investigar.

\section{Objetivos}

\subsection{Objetivo Geral}

Síntese do que se pretende alcançar com a pesquisa.

\subsection{Objetivos Específicos}

\begin{itemize}
  \item Objetivo específico 1
  \item Objetivo específico 2
  \item Objetivo específico 3
\end{itemize}

% --- 2. REFERENCIAL TEÓRICO ---
\chapter[Referencial Teórico]{Referencial Teórico}

\section{Tema Central}

Fundamentação teórica do constructo principal.

\subsection{Subtema 1}

Desenvolvimento com citações:

\begin{citacao}
A citação direta longa deve ser formatada em parágrafo independente, recuado
4 cm da margem esquerda, em fonte 10, espaçamento simples e sem aspas,
conforme a NBR 10520:2023. A referência (Autor, ano, p. X) deve ser indicada
fora do bloco.
\end{citacao}

\subsection{Subtema 2}

Desenvolvimento do segundo constructo.

\section{Estado da Arte}

Revisão sistemática ou integrativa da literatura recente sobre o tema,
demonstrando o que já foi pesquisado e quais lacunas persistem.

% --- 3. METODOLOGIA ---
\chapter[Metodologia]{Metodologia}

\section{Classificação da Pesquisa}

\section{Contexto e Participantes}

\section{Instrumentos e Procedimentos}

\section{Procedimentos de Análise}

\section{Aspectos Éticos}

% --- 4. RESULTADOS ---
\chapter[Resultados]{Resultados}

\section{Análise Descritiva}

\section{Análise Inferencial}

\begin{table}[H]
\centering
\caption{Resultados da análise estatística.}
\label{tab:resultados}
\begin{tabular}{lccc}
\toprule
\textbf{Variável} & \textbf{Pré-teste} & \textbf{Pós-teste} & \textbf{p} \\
\midrule
Indicador A & 45,2 $\pm$ 3,1 & 52,8 $\pm$ 2,9 & $< 0,05$ \\
Indicador B & 78,6 $\pm$ 4,2 & 81,3 $\pm$ 3,8 & $0,12$ \\
\bottomrule
\end{tabular}
\fonte{Elaborado pelo autor (2026).}
\end{table}

% --- 5. DISCUSSÃO ---
\chapter[Discussão]{Discussão}

Interpretação dos resultados à luz da literatura. Comparação com estudos
prévios, limitações e implicações.

% --- 6. CONCLUSÃO ---
\chapter[Conclusão]{Conclusão}

Síntese dos achados, contribuições, limitações e sugestões para trabalhos
futuros.

% ============================================================================
% ELEMENTOS PÓS-TEXTUAIS
% ============================================================================
\postextual

% --- REFERÊNCIAS ---
\bibliography{dissertacao_modelo_abnt}

% --- GLOSSÁRIO (opcional) ---
% \begin{glossario}
%   \item[Termo 1] Definição do termo 1.
%   \item[Termo 2] Definição do termo 2.
% \end{glossario}

% --- APÊNDICES (opcionais) ---
\begin{apendicesenv}
\partapendices

\chapter{Instrumento de Coleta de Dados}
Texto do instrumento (questionário, roteiro de entrevista, etc.)

\end{apendicesenv}

% --- ANEXOS (opcionais) ---
% \begin{anexosenv}
% \partanexos
%
% \chapter{Documento Original}
% Texto do documento anexado.
%
% \end{anexosenv}

% --- ÍNDICE REMISSIVO (opcional) ---
% \printindex

\end{document}
